\chapter{Floods and Cyclones: Health Risks}

\section*{Definition and Overview}
Floods and cyclones are natural disasters that strike parts of Australia regularly, and they bring a range of health risks that continue long after the waters recede and the winds stop. The immediate dangers are drowning, injury, and hypothermia, while the aftermath brings contaminated floodwater, infectious disease, hazardous debris, carbon monoxide poisoning from generators, mental health impacts, and problems with clean water, food, and medicines. People with chronic conditions are particularly vulnerable when services and supplies are disrupted. Understanding these risks and preparing for them is the most effective way to protect health. This chapter outlines the health risks of floods and cyclones and how to manage them.

\section*{Epidemiology}
Australia experiences some of the most variable flood and cyclone activity in the world, with cyclones affecting the north and floods affecting many regions, including major events in Queensland, New South Wales, and Victoria. Floods and cyclones cause more injury, disease, and mental health burden in Australia than most other natural disasters. Drowning and trauma are the leading causes of death in the acute phase, while the aftermath is followed by surges in gastrointestinal illness, skin infections, leptospirosis, mould-related respiratory problems, and mental health conditions including post-traumatic stress, depression, and anxiety that can last for years. Vulnerable groups, including older adults, people with disability, and those with chronic illness, carry a disproportionate burden.

\section*{Aetiology and Pathophysiology}
\begin{itemize}
    \item \textbf{Floodwater:} contains sewage, chemicals, animal waste, and debris; contact or ingestion causes gastroenteritis, skin and wound infections, and diseases such as leptospirosis and melioidosis in affected regions
    \item \textbf{Drowning:} the leading cause of death, including people attempting to drive or walk through floodwater
    \item \textbf{Injury:} trauma from debris, unstable buildings, electrocution from damaged power lines, and chainsaw and generator accidents during clean-up
    \item \textbf{Carbon monoxide:} generators, heaters, and cooking equipment used indoors or in confined spaces cause poisoning, which can be fatal
    \item \textbf{Mould:} damp homes grow mould, worsening asthma, allergies, and other respiratory conditions
    \item \textbf{Contaminated water and food:} damage to water supplies and refrigeration causes waterborne and foodborne illness
    \item \textbf{Mosquito-borne disease:} standing water breeds mosquitoes, raising the risk of Ross River virus, Barmah Forest virus, and dengue in the north
    \item \textbf{Disruption of care:} loss of power, medicines, and access to health services harms people with chronic conditions
    \item \textbf{Mental health:} the trauma of the event, loss of home and possessions, and prolonged recovery cause anxiety, depression, and post-traumatic stress
    \item \textbf{Heat and hypothermia:} exposure during and after the event, depending on season
\end{itemize}

\section*{Clinical Features}
Health effects appear in phases after a flood or cyclone.
\begin{itemize}
    \item Immediate phase: drowning, injuries, hypothermia, and electrocution
    \item Early aftermath (days to weeks): gastroenteritis from contaminated water, wound infections, leptospirosis with fever and muscle pain, and carbon monoxide poisoning from generators
    \item Later aftermath (weeks to months): mould-related respiratory illness, mosquito-borne disease, skin conditions, and mental health problems
    \item In people with chronic illness: worsening of asthma, diabetes, heart disease, and kidney disease from disrupted treatment
    \item In children and older adults: dehydration, infection, and distress
    \item Psychological features: insomnia, anxiety, irritability, difficulty concentrating, and avoidance
    \item The risk of tetanus and other wound infections during clean-up
\end{itemize}

\section*{Differential Diagnosis}
Illness after a flood needs sorting into its possible causes.
\begin{itemize}
    \item Gastrointestinal illness from contaminated water versus food poisoning
    \item Fever and muscle pain from leptospirosis, melioidosis, or mosquito-borne viruses
    \item Respiratory symptoms from mould, asthma, or infection
    \item Carbon monoxide poisoning, which mimics flu with headache, nausea, and confusion
    \item Mental health conditions, which overlap with the normal stress response
    \item The need to consider these in anyone unwell after a flood, even without a clear history
\end{itemize}

\section*{When to Seek Medical Care}
Anyone with fever, vomiting, diarrhoea, worsening respiratory symptoms, infected wounds, or significant distress after a flood or cyclone should seek medical care. People with chronic conditions should re-establish care early and restock medicines.

\textbf{Call 000 (or local emergency services) for:}
\begin{itemize}
    \item Drowning, unconsciousness, or difficulty breathing
    \item Severe bleeding, trauma, or suspected spinal injury
    \item Carbon monoxide exposure: headache, nausea, confusion, or collapse in a building with a generator or heater running
    \item Chest pain, severe breathlessness, or signs of stroke or heart attack
    \item Suicidal thoughts or behaviour
    \item Signs of sepsis: fever with confusion, rapid heart rate, or collapse
\end{itemize}

\section*{Diagnosis and Investigations}
\begin{itemize}
    \item \textbf{History:} exposure to floodwater, injuries, symptoms, and disruption to medicines and care
    \item \textbf{Physical examination:} assessment of wounds, respiratory system, and hydration
    \item \textbf{Blood and stool tests:} for leptospirosis, melioidosis, gastrointestinal infections, and mosquito-borne viruses when suspected
    \item \textbf{Carbon monoxide testing:} for people exposed to generators or combustion in enclosed spaces
    \item \textbf{Respiratory review:} for people with asthma and mould exposure
    \item \textbf{Mental health assessment:} screening for anxiety, depression, and post-traumatic stress
    \item \textbf{Review of chronic disease:} blood pressure, blood sugar, and kidney function checks after disruption of treatment
\end{itemize}

\section*{Management}
\begin{itemize}
    \item \textbf{Acute emergencies:} rescue, first aid, and urgent hospital care for drowning, trauma, and carbon monoxide poisoning
    \item \textbf{Water safety:} treat all floodwater as contaminated; boil or use bottled water until supply is confirmed safe
    \item \textbf{Wound care:} clean wounds thoroughly, keep them covered, and seek care for deep or infected wounds, with tetanus vaccination updated
    \item \textbf{Carbon monoxide prevention:} never run generators, heaters, or engines indoors or near open windows
    \item \textbf{Mould management:} dry homes quickly, remove soaked materials, and use masks during clean-up; treat worsening asthma promptly
    \item \textbf{Infectious disease:} treat gastroenteritis with fluids and rehydration, and specific infections with antibiotics or antivirals under medical care
    \item \textbf{Mosquito protection:} repellent, protective clothing, and removal of standing water
    \item \textbf{Chronic disease:} restock medicines, re-establish care with the GP, and seek review of conditions affected by disruption
    \item \textbf{Mental health:} psychological first aid, counselling, and support services, with urgent care for crisis
    \item \textbf{Community recovery:} access recovery services, financial assistance, and support groups
\end{itemize}

\section*{Complications}
\begin{itemize}
    \item Death and disability from drowning, trauma, and carbon monoxide poisoning
    \item Outbreaks of gastroenteritis and serious infections in affected communities
    \item Worsening of chronic disease with poor outcomes
    \item Mould-related respiratory disease and new-onset asthma
    \item Post-traumatic stress, depression, anxiety, and increased suicide risk
    \item Long-term financial and housing stress that compounds health problems
\end{itemize}

\section*{Prognosis and Outcomes}
Most health effects of floods and cyclones are preventable or treatable with prompt care, and communities recover with time and support. Early treatment of infections, safe water and food handling, and mould control prevent most of the aftermath illness. Mental health outcomes improve substantially with timely support, and people with chronic conditions do well when their care is re-established quickly. Preparation is the most powerful factor: households with plans, supplies, and knowledge of the risks emerge healthier and recover faster.

\section*{Prevention and Self-Care}
\begin{itemize}
    \item Know your flood and cyclone risk, and have an emergency plan and kit with medicines, water, food, and documents
    \item Never drive, walk, or ride through floodwater; more deaths occur in vehicles than anywhere else
    \item Treat all floodwater as contaminated and protect skin from contact
    \item Use generators and fuel-burning appliances only outdoors and away from windows
    \item Boil or treat water until authorities confirm it is safe
    \item Dry the home quickly after flooding and wear a mask when cleaning mould
    \item Restock medicines and re-establish care early if services were disrupted
    \item Seek mental health support early, and look out for family and neighbours
\end{itemize}

\section*{Sources and Further Reading}
The following reputable sources provide further information for healthcare students and patients seeking reliable, current advice about the health risks of floods and cyclones.
\begin{itemize}
    \item Healthdirect Australia. \textit{Flood safety and health.} https://www.healthdirect.gov.au/
    \item Australian Institute of Health and Welfare. \textit{Health impacts of natural disasters.} https://www.aihw.gov.au/
    \item Emergency Management Australia. \textit{Preparing for floods and cyclones.} https://www.ema.gov.au/
    \item Queensland Health. \textit{Flood health information.} https://www.health.qld.gov.au/
    \item Centers for Disease Control and Prevention. \textit{Health and safety after floods.} https://www.cdc.gov/
    \item World Health Organization. \textit{Natural disasters and health.} https://www.who.int/
\end{itemize}
